\chapter{Angiotensin-Converting Enzyme (ACE) Inhibitors}

\section*{Definition and Overview}
Angiotensin-converting enzyme (ACE) inhibitors are a class of medicines that relax blood vessels and lower blood pressure. They are used mainly to treat hypertension (high blood pressure) and heart failure, and to protect the kidneys in people with diabetes and chronic kidney disease. Commonly used examples include ramipril, perindopril, enalapril, lisinopril, and captopril.

ACE inhibitors work by blocking angiotensin-converting enzyme, the enzyme that converts angiotensin I into angiotensin II. Angiotensin II is a potent substance that narrows blood vessels, stimulates the release of aldosterone, and promotes salt and fluid retention. By reducing its production, ACE inhibitors lower blood pressure, reduce the workload on the heart, and decrease pressure within the kidney's filtering units. They are among the most extensively studied medicines in clinical practice and are recommended as first-line treatment in many guidelines because they reduce cardiovascular events and improve survival in heart failure.

\section*{Epidemiology}
High blood pressure affects roughly one in three adults worldwide and is a leading cause of heart attack, stroke, heart failure, kidney disease, and premature death. ACE inhibitors are among the most widely prescribed medicines for hypertension and heart failure, particularly in older adults, in people with diabetes, and in those who have had a heart attack or stroke.

In many national and international guidelines, ACE inhibitors are a preferred first-line option for hypertension, especially when other conditions such as diabetes, kidney disease, or heart failure are present. Their use has increased substantially over recent decades, and they remain a cornerstone of cardiovascular prevention across all age groups in whom they are not contraindicated.

\section*{Aetiology and Pathophysiology}
The renin--angiotensin--aldosterone system is a hormonal cascade that regulates blood pressure and fluid balance. When blood pressure falls or sodium is low, the kidney releases renin, which converts the liver protein angiotensinogen into angiotensin I. Angiotensin-converting enzyme in the lungs and other tissues then converts angiotensin I into angiotensin II.

Angiotensin II exerts three major effects relevant to disease: it is a powerful constrictor of blood vessels, raising peripheral resistance and blood pressure; it stimulates the adrenal glands to release aldosterone, which makes the kidney retain salt and water; and it promotes structural changes (remodelling) in the heart and blood vessel walls. ACE inhibitors interrupt this system:

\begin{itemize}
    \item \textbf{Blocking angiotensin II production:} blood vessels widen and blood pressure falls
    \item \textbf{Reducing aldosterone:} the body excretes more salt and water, reducing fluid overload
    \item \textbf{Lowering glomerular filtration pressure:} the pressure inside the kidney's filters falls, which protects the kidneys, particularly in diabetic nephropathy
    \item \textbf{Reducing cardiac strain:} afterload falls and heart-failure remodelling is slowed
    \item \textbf{Preserving bradykinin:} ACE normally breaks down bradykinin, a vasodilator; inhibiting this raises bradykinin levels, which contributes to both benefit and the dry-cough side effect
\end{itemize}

\section*{Clinical Features}
Most people taking an ACE inhibitor feel no direct symptoms, and the benefit is measured by blood pressure, kidney function, and heart symptoms. When adverse effects occur they typically include:

\begin{itemize}
    \item \textbf{Dry cough:} a persistent, tickly, non-productive cough, the most common side effect, affecting roughly 5--10\% of people
    \item \textbf{Dizziness or light-headedness:} especially after the first dose, on standing, or with dehydration
    \item \textbf{Hyperkalaemia:} a rise in blood potassium detected on blood tests
    \item \textbf{Rising creatinine:} an early, generally acceptable dip in kidney function followed by stabilisation
    \item \textbf{Altered taste, headache, or fatigue:} uncommon and usually transient
    \item \textbf{Rash or itching:} a possible hypersensitivity reaction
    \item \textbf{Angioedema:} sudden swelling of the lips, face, tongue, or throat; rare but potentially life-threatening
\end{itemize}

\section*{Differential Diagnosis}
New symptoms while taking an ACE inhibitor may be due to the medicine or to an unrelated condition:

\begin{itemize}
    \item Cough from the medicine versus a chest infection, asthma, reflux, or post-nasal drip
    \item Dizziness from low blood pressure versus dehydration, anaemia, or an abnormal heart rhythm
    \item Angioedema from the medicine versus an allergic reaction to food, insect stings, or another drug
    \item Worsening kidney function from the medicine versus progression of underlying kidney disease
    \item Raised potassium from the medicine versus renal failure, or interaction with potassium-sparing diuretics or NSAIDs
\end{itemize}

The clinical context, timing of symptoms in relation to starting the drug, and the response to dose adjustment usually distinguish these possibilities.

\section*{When to Seek Medical Care}
Tell your doctor promptly if you develop a new persistent cough, feel persistently dizzy or faint, or are prescribed other medicines that interact, such as non-steroidal anti-inflammatory drugs (NSAIDs), diuretics, or potassium supplements. People who are pregnant, planning a pregnancy, or breastfeeding should discuss their medicines urgently.

\textbf{Call 000 (or local emergency services) for:}
\begin{itemize}
    \item Swelling of the lips, face, tongue, or throat, or difficulty breathing or swallowing (angioedema)
    \item Fainting, collapse, or severe dizziness
    \item Chest pain, palpitations, or a very fast or irregular heartbeat
\end{itemize}

\section*{Diagnosis and Investigations}
Monitoring is essential to use ACE inhibitors safely:

\begin{itemize}
    \item \textbf{Blood pressure measurement:} before starting and at each review to assess response
    \item \textbf{Kidney function and electrolytes:} creatinine, urea, and potassium are measured before starting, within one to two weeks after any dose change, and periodically thereafter
    \item \textbf{Regular review:} kidney function and potassium are usually rechecked at set intervals, more frequently in older people and those with kidney disease
    \item \textbf{Pregnancy test:} considered before starting in women of childbearing potential, since ACE inhibitors are harmful in pregnancy
    \item \textbf{Medication review:} to detect interactions with NSAIDs, diuretics, potassium-sparing agents, and other blood pressure drugs
\end{itemize}

\section*{Management}
\begin{itemize}
    \item \textbf{Start low and go slow:} the dose is increased gradually to avoid first-dose hypotension
    \item \textbf{Consistency:} take the medicine at the same time each day, usually once daily
    \item \textbf{First-dose caution:} take the first dose when you can sit or lie down if you feel dizzy
    \item \textbf{Troublesome cough:} switching to an angiotensin receptor blocker (ARB) resolves the cough in most people
    \item \textbf{Angioedema:} the ACE inhibitor must be stopped permanently; an ARB is used only with caution and specialist advice
    \item \textbf{Combination therapy:} often combined with a calcium channel blocker or diuretic to reach blood pressure targets
    \item \textbf{Contraindicated in pregnancy:} treatment is changed to a pregnancy-safe alternative before conception
    \item \textbf{Adjunct lifestyle advice:} reduce salt, limit alcohol, exercise, and maintain a healthy weight alongside the medicine
\end{itemize}

\section*{Complications}
\begin{itemize}
    \item Dry cough, which can lead to poor adherence if not recognised and managed
    \item Hyperkalaemia, which can cause dangerous heart rhythms in severe cases
    \item Worsening renal function, particularly with dehydration, bilateral renal artery stenosis, or concurrent NSAID use
    \item Angioedema, which can threaten the airway and requires emergency treatment
    \item Hypotension and falls, especially in older people and after the first dose
    \item Fetal harm if used during pregnancy, including renal and other developmental problems
\end{itemize}

\section*{Prognosis and Outcomes}
ACE inhibitors are among the best-evidenced medicines in clinical medicine. In heart failure they reduce hospital admissions and improve survival. In people with diabetes and proteinuric kidney disease they slow the decline in kidney function and reduce the risk of progressing to end-stage renal disease. In hypertension they lower the risk of stroke, heart attack, and heart failure, and after a myocardial infarction they improve survival. Provided kidney function and potassium are monitored and the medicine is taken consistently, most people tolerate ACE inhibitors well and derive benefit for many years.

\section*{Prevention and Self-Care}
\begin{itemize}
    \item Do not stop the medicine suddenly, since blood pressure can rebound
    \item Keep to the scheduled blood tests for kidney function and potassium
    \item Stay well hydrated, but avoid over-the-counter NSAIDs and potassium-rich salt substitutes without medical advice
    \item Rise slowly from sitting or lying to reduce dizziness and the risk of falls
    \item Keep an up-to-date medicine list and share it with every health professional you see
    \item If pregnant, planning pregnancy, or breastfeeding, discuss your medicines with your doctor before making any change
    \item Report any swelling of the face, lips, or tongue immediately
\end{itemize}

\section*{Sources and Further Reading}
The following sources provide reliable further information on ACE inhibitors:

\begin{itemize}
    \item NHS. \textit{Information on ACE inhibitors, including how they work and common side effects.} https://www.nhs.uk/conditions/ace-inhibitors/
    \item MedlinePlus. \textit{Overview of ACE inhibitor medicines, their uses, and precautions.} https://medlineplus.gov/aceinhibitors.html
    \item Mayo Clinic. \textit{ACE inhibitors: description, benefits, and possible side effects.} https://www.mayoclinic.org/diseases-conditions/high-blood-pressure/in-depth/ace-inhibitors/art-20048180
    \item NICE. \textit{Clinical guideline on hypertension in adults, including first-line treatment recommendations.} https://www.nice.org.uk/guidance/ng136
    \item MSD Manual. \textit{Angiotensin-converting enzyme inhibitors in the professional health section.} https://www.msdmanuals.com/professional/cardiovascular-disorders/antihypertensive-drugs/ace-inhibitors
\end{itemize}
